\documentclass[11pt]{article}

\usepackage[utf8]{inputenc}
\usepackage[T1]{fontenc}
\usepackage[spanish,es-tabla]{babel}
\usepackage{lmodern}
\usepackage[a4paper,margin=2.2cm]{geometry}
\usepackage{graphicx}
\usepackage{array}
\usepackage{longtable}
\usepackage{booktabs}
\usepackage[table]{xcolor}
\usepackage{enumitem}
\usepackage[hidelinks]{hyperref}

\graphicspath{{capturas/}}
\setlength{\parskip}{0.5em}
\setlength{\parindent}{0pt}
\renewcommand{\arraystretch}{1.35}
\setlist{nosep,leftmargin=1.2em}

\definecolor{marca}{HTML}{0B3A5C}
\definecolor{borde}{HTML}{B8C4CE}
\definecolor{suave}{HTML}{F3F6F9}

% Captura: si el archivo existe en capturas/ se muestra; si no, queda el espacio con su nombre.
\newcommand{\captura}[1]{%
  \IfFileExists{capturas/#1.png}{\includegraphics[width=\linewidth,height=5.2cm,keepaspectratio]{#1.png}}{%
  \IfFileExists{capturas/#1.jpg}{\includegraphics[width=\linewidth,height=5.2cm,keepaspectratio]{#1.jpg}}{%
  {\setlength{\fboxsep}{0pt}\fcolorbox{borde}{suave}{\parbox[c][3cm][c]{\dimexpr\linewidth-2\fboxrule\relax}{\centering\small\color{gray}\texttt{capturas/#1.png}}}}}}}

% Un problema completo con sus fotos.
% #1 número (h1, h2...), #2 título, #3 quién lo encontró, #4 prioridad, #5 qué pasó, #6 qué se cambió
\newcommand{\problema}[6]{%
  \par\bigskip\noindent\begin{minipage}{\linewidth}
  {\large\textbf{\color{marca}#1}\ \ \textbf{#2}}\par\smallskip
  {\small\colorbox{suave}{\strut\ \textbf{Quién lo encontró:} #3 \qquad \textbf{Prioridad:} #4\ }}\par\smallskip
  \textbf{Qué pasó.} #5\par\smallskip
  \textbf{Qué se cambió.} #6\par\medskip
  \begin{minipage}[t]{0.485\linewidth}\centering
    {\small\textbf{Antes}}\par\smallskip\captura{#1-antes}
  \end{minipage}\hfill
  \begin{minipage}[t]{0.485\linewidth}\centering
    {\small\textbf{Después}}\par\smallskip\captura{#1-despues}
  \end{minipage}
  \end{minipage}\par}

% Paso del ciclo de pruebas.
\newcommand{\paso}[1]{{\setlength{\fboxsep}{5pt}\fcolorbox{borde}{suave}{\textbf{\color{marca}#1}}}}

\title{Reporte de validación RITE\\[0.3em]
\large Loard - Hito 2 - Módulo III: QC de usabilidad y accesibilidad}
\author{}
\date{16 de noviembre de 2026}

\begin{document}
\maketitle

% =====================================================================
\section{Introducción}

Este reporte recoge las pruebas del prototipo de Loard con usuarios reales.
Se usa el método RITE: se prueba, se corrige lo que falla y se vuelve a probar.
El objetivo es comprobar que las personas hacen las tareas principales sin ayuda
y dejar evidencia de cada cambio.

% =====================================================================
\section{Cómo se hizo}

\begin{center}
\paso{Ronda 1} \ $\rightarrow$\ \paso{Ronda 2} \ $\rightarrow$\ \paso{Ajustes al prototipo} \ $\rightarrow$\ \paso{Ronda 3}
\end{center}

\begin{enumerate}[leftmargin=1.5em]
  \item \textbf{Quiénes.} Tres personas reales que usarían Loard en su trabajo: dos
        residentes de obra y un director de proyectos.
  \item \textbf{Sesión.} Una persona a la vez, con el prototipo abierto.
  \item \textbf{Pensar en voz alta.} Mientras hacía cada tarea, la persona decía lo que veía,
        lo que buscaba y lo que esperaba que pasara.
  \item \textbf{Instrucciones sin pistas.} Cada tarea decía qué había que lograr, no cómo:
        no se nombraron botones, menús ni pantallas. Quien dirigía la prueba ayudó lo mínimo,
        solo para que la prueba fluyera.
  \item \textbf{Revisión libre.} Los aportes no se limitaron a los problemas de lo que hay
        hoy en el prototipo. Después de las tareas se usó el espacio para sacar más ideas:
        la persona recorrió el prototipo a su gusto, explicó cómo es el trabajo diario en
        la obra y propuso mejoras.
  \item \textbf{Encuesta.} Al terminar se respondió la encuesta SUS (sección~\ref{sec:sus}).
  \item \textbf{Rondas.} Se hicieron las rondas 1 y 2, se ajustó el prototipo con lo
        encontrado y después se hizo la ronda 3 con las mismas tareas de la ronda 1.
\end{enumerate}

% =====================================================================
\section{Participantes}

\begin{longtable}{>{\raggedright\arraybackslash}p{2.5cm} >{\raggedright\arraybackslash}p{4.1cm} >{\raggedright\arraybackslash}p{5cm} >{\raggedright\arraybackslash}p{2.6cm}}
\toprule
\textbf{Participante} & \textbf{Rol} & \textbf{Experiencia} & \textbf{Fecha} \\
\midrule
\endhead
\bottomrule
\endfoot
Residente 1 & Residente de obra & Usa la suite de Microsoft para todo lo relacionado con su trabajo. & 16 de noviembre de 2026 \\
Director & Director de proyectos. También conoce una parte de SST & Usa la suite de Microsoft para todo lo relacionado con su trabajo. & 16 de noviembre de 2026 \\
Residente 2 & Residente de obra & Usa la suite de Microsoft para todo lo relacionado con su trabajo. & 16 de noviembre de 2026 \\
\end{longtable}

% =====================================================================
\section{Tareas}

\begin{longtable}{>{\raggedright\arraybackslash}p{1.15cm} >{\raggedright\arraybackslash}p{1.7cm} >{\raggedright\arraybackslash}p{9.35cm} >{\raggedright\arraybackslash}p{2.1cm}}
\toprule
\textbf{Tarea} & \textbf{HU} & \textbf{Instrucción para el usuario} & \textbf{Quién la hizo} \\
\midrule
\endhead
\bottomrule
\endfoot
\textbf{T1} &
\mbox{HU-M02}\newline \mbox{HU-M10} &
``Hoy su cuadrilla terminó un trabajo en el frente de obra. Registre el corte del día
con las cantidades y el personal que trabajó. Incluya también a un trabajador que
está bloqueado por SST.'' &
Residente 1 y Residente 2 \\

\textbf{T2} &
\mbox{HU-M04} &
``Usted revisó un corte y las cantidades no coinciden con lo que
vio en la obra. Haga que el residente lo corrija y déjele claro qué debe cambiar.'' &
Director \\

\textbf{T3} &
\mbox{HU-M06} &
``Le entregaron un soporte que pide el dossier del proyecto. Súbalo en
el lugar que le corresponde.'' &
Residente 1 y Residente 2 \\

\textbf{T4} &
\mbox{HU-M12} &
``Se necesita saber cómo va el proyecto. Averigüe cuánto avance lleva y si hay algo
atrasado o pendiente.'' &
Director \\
\end{longtable}

% =====================================================================
\section{Encuesta SUS}\label{sec:sus}

Al terminar, cada persona marcó de 1 a 5 qué tan de acuerdo estaba con cada frase:
1 = muy en desacuerdo, 5 = muy de acuerdo.

\begin{longtable}{>{\raggedright\arraybackslash}p{0.75cm} >{\raggedright\arraybackslash}p{9.9cm} >{\centering\arraybackslash}p{1.9cm} >{\centering\arraybackslash}p{1.9cm}}
\toprule
\textbf{N.º} & \textbf{Pregunta} & \textbf{Residente de obra} & \textbf{Director de proyectos} \\
\midrule
\endhead
\bottomrule
\endfoot
1  & Creo que me gustaría utilizar este sistema con frecuencia. & 5 & 5 \\
2  & Encontré el sistema innecesariamente complejo. & 1 & 2 \\
3  & Pensé que el sistema era fácil de usar. & 5 & 4 \\
4  & Creo que necesitaría el apoyo de un técnico para poder utilizar este sistema. & 1 & 2 \\
5  & Encontré que las diversas funciones de este sistema estaban bien integradas. & 3 & 3 \\
6  & Pensé que había demasiada inconsistencia en el sistema. & 1 & 1 \\
7  & Imagino que la mayoría de las personas aprenderían a utilizar este sistema muy rápidamente. & 5 & 4 \\
8  & Encontré el sistema muy incómodo de utilizar. & 1 & 1 \\
9  & Me sentí muy seguro utilizando el sistema. & 4 & 5 \\
10 & Necesité aprender muchas cosas antes de poder empezar a utilizar este sistema. & 1 & 2 \\
\midrule
   & \textbf{Puntaje SUS (0 a 100)} & \textbf{92,5} & \textbf{82,5} \\
\end{longtable}

\textbf{Cómo se calcula.} En las preguntas impares se resta 1 a la respuesta; en las pares
se resta la respuesta a 5. Se suman los 10 resultados y se multiplica por 2,5:
\[
\mathrm{SUS} = 2{,}5 \times \Big[\,(r_1-1) + (5-r_2) + (r_3-1) + \cdots + (5-r_{10})\,\Big]
\]

\textbf{Promedio SUS:} 87,5 \qquad \textbf{Meta:} 70 o más \qquad \textbf{¿Se cumple?} Sí

% =====================================================================
\section{Evidencia: tabla de ``Problema detectado'' vs. ``Cambio aplicado en el prototipo''}

En las pruebas salieron siete problemas de uso. La tabla los resume y, después, cada uno
se explica con quién lo encontró, qué pasó, qué se cambió en el prototipo y la pantalla
antes y después del cambio.

\textbf{Prioridad:} P0 impide terminar la tarea; P1 la tarea se logra con mucha dificultad;
P2 molesta o demora, pero se logra; P3 detalle menor.

\begin{longtable}{>{\raggedright\arraybackslash}p{0.6cm} >{\raggedright\arraybackslash}p{4.6cm} >{\raggedright\arraybackslash}p{4.7cm} >{\raggedright\arraybackslash}p{2.3cm} >{\centering\arraybackslash}p{1.95cm}}
\toprule
\textbf{\#} & \textbf{Problema detectado} & \textbf{Cambio aplicado en el prototipo} & \textbf{Quién lo encontró} & \textbf{Prioridad} \\
\midrule
\endhead
\bottomrule
\endfoot
h1 & No se podían subir soportes al registrar el corte. & Los soportes se agregan en el mismo formulario y quedan en el corte al guardar. & Residente 1 & P0 \\
h2 & No se podía escribir una actividad que no estuviera en la lista. & Opción ``Otra actividad (escribir)'' con nombre, unidad y cantidad. & Residente 1 & P1 \\
h3 & La cantidad no subía de uno en uno. & La cantidad sube de uno en uno y solo acepta números enteros. & Residente 1 & P2 \\
h4 & La búsqueda de cortes solo miraba el estado elegido. & La búsqueda revisa todos los cortes y hay filtros por fecha, ciudad y cliente. & Director & P1 \\
h5 & Ciudad, residente y contrato se veían pequeños en el corte. & Los tres datos van en letra grande arriba del corte. & Director & P2 \\
h6 & No se reconocía el botón de cerrar sesión. & El botón usa el símbolo de apagar. & Director & P3 \\
h7 & Los cuadros de actividad y cantidad no estaban alineados. & Los cuadros quedan a la misma altura y con el mismo tamaño. & Residente 2 & P3 \\
\end{longtable}

\subsection*{Cada problema con su antes y después}

\problema{h1}{No se podían subir soportes al registrar el corte}{Residente 1}{P0}
{Al registrar el corte, el residente quiso adjuntar las fotos y el PDF del trabajo. La sección
``Soportes del corte'' solo mostraba un texto que pedía guardar primero el corte, y no encontró
dónde subirlos ni los vio después. No pudo terminar esa parte de la tarea.}
{Ahora los soportes se agregan en el mismo formulario del corte: se arrastran o se eligen
los archivos (fotos o PDF), cada uno aparece en la lista apenas se sube y se puede quitar.
Al guardar el corte, los soportes quedan dentro de él.}

\problema{h2}{No se podía escribir una actividad propia}{Residente 1}{P1}
{El residente había hecho un trabajo que no estaba en la lista (pintura de tubería). La lista
solo traía actividades ya cargadas y no había forma de escribir otra, así que no pudo
registrar lo que de verdad se hizo.}
{La lista tiene la opción ``Otra actividad (escribir)''. Al elegirla se escribe el nombre de la
actividad y se escoge su unidad de medida (m, m\textsuperscript{2}, m\textsuperscript{3}, kg
o unidades), junto con la cantidad.}

\problema{h3}{La cantidad no subía de uno en uno}{Residente 1}{P2}
{Al anotar la cantidad ejecutada, el cuadro aceptaba decimales y no subía de uno en uno.
Costaba dejar el número exacto y era fácil equivocarse.}
{La cantidad ejecutada sube y baja de uno en uno (1, 2, 3\ldots) y solo acepta números
enteros. Si se escribe un decimal, el formulario avisa.}

\problema{h4}{La búsqueda de cortes dependía del estado}{Director}{P1}
{El director buscó un corte por su nombre y no apareció, porque la lista solo mostraba los
cortes del estado que estaba elegido (por ejemplo, ``En revisión''). Tampoco pudo buscar por
fecha, ciudad ni cliente.}
{Al escribir en el buscador se revisan todos los cortes, sin importar su estado, y la pantalla
lo avisa. Se agregaron filtros por fecha (desde y hasta), ciudad y cliente, y la lista muestra
la ciudad y el cliente de cada corte.}

\problema{h5}{Ciudad, residente y contrato se veían pequeños}{Director}{P2}
{Al abrir un corte, al director le costó ubicar de qué ciudad era, qué residente lo hizo y a
qué contrato u obra pertenecía. Esos datos estaban en letra pequeña, mezclados con otros.}
{Arriba del corte hay un recuadro con los tres datos en letra grande: ciudad, residente y
contrato u obra, con el cliente debajo.}

\problema{h6}{No se reconocía el botón de cerrar sesión}{Director}{P3}
{El director no reconoció el botón para cerrar sesión. Tenía una flecha de salida y él
esperaba el símbolo de apagar, que es el que usa todos los días en su correo.}
{El botón de cerrar sesión ahora usa el símbolo de apagar.}

\problema{h7}{Los cuadros de actividad y cantidad no estaban alineados}{Residente 2}{P3}
{En ``Actividades ejecutadas'', el cuadro de la actividad y el de la cantidad quedaban a
distinta altura y con distinto tamaño. Se veía desordenado.}
{Los dos cuadros y el botón de quitar quedan a la misma altura y con el mismo tamaño.
La unidad de medida pasó al título del cuadro de cantidad.}

% =====================================================================
\section{Pedidos para próximas versiones}

Además de los problemas, las personas pidieron funciones nuevas. No son fallas del
prototipo: son necesidades del trabajo diario que quedan anotadas para más adelante.

\textbf{Residentes de obra}
\begin{itemize}
  \item \textbf{Actividades más detalladas.} En tubería, separar los metros tendidos de la cantidad de soldaduras y anotar el diámetro (4, 6 u 8 pulgadas), porque no todos rinden igual. En obra civil, agregar los pasos previos (excavación, formaleta, relleno, compactación, solado).
  \item \textbf{Soldadores.} Registrar su calificación vigente en soldadura, además de los cursos.
  \item \textbf{Materiales.} Registrar el lote que llega al proyecto (por ejemplo, 4 cajas de soldadura) en lugar de lo que se gasta cada día.
  \item \textbf{Herramientas y equipos.} Anotar cuáles se usaron en el día y llevar su mantenimiento igual que con las personas: último mantenimiento y horas que faltan para el siguiente, con aviso cuando falte poco.
  \item \textbf{Clima.} Anotar el clima del día (sol, horas de lluvia) para explicar el avance.
  \item \textbf{Cortes.} El corte suele ser mensual. Debe llegarle al director o al gerente y poder descargarse para enviarlo al cliente.
  \item \textbf{Cada quien ve lo suyo.} Que cada residente vea solo los cortes y documentos de su obra o ciudad.
  \item \textbf{Proyectos nuevos.} Un botón para registrar un contrato apenas se apruebe, con aviso al director.
  \item \textbf{Avisos.} Correo cuando devuelvan un corte o haya que revisar el dossier, y recordatorio de los cortes cuando se acerque la facturación.
  \item \textbf{Más adelante.} Pasar las horas trabajadas a la nómina, llevar la caja y el balance de cada obra, y un módulo de cotizaciones con precios unitarios que el director revise en línea y se exporte a PDF o Excel.
\end{itemize}

\textbf{Director de proyectos}
\begin{itemize}
  \item \textbf{Elegir la obra.} Escoger en el tablero de qué obra se quieren ver los datos.
  \item \textbf{Avance.} Medirlo comparando los días cotizados con los días que van, junto con las cantidades ejecutadas.
  \item \textbf{Revisión diaria.} Una lista diaria para comprobar que equipos y herramientas están en condiciones de trabajar.
  \item \textbf{Equipos y calibraciones.} Llevar en un solo lugar el mantenimiento y los certificados de los equipos que viajan entre obras; hoy están repartidos en correos.
  \item \textbf{Más adelante.} Ver en el tablero el dinero ejecutado y lo facturado, conectar el plan de trabajo de MS Project y que el encargado de SST de cada obra actualice sus propios cursos y certificados.
\end{itemize}

% =====================================================================
\section{Conclusión}

Las tres personas entendieron el prototipo y lo recorrieron con muy poca ayuda. La encuesta SUS dio
92,5 y 82,5 puntos, con un promedio de 87,5, por encima de la meta de 70.

Se encontraron seis problemas de uso. El más importante fue que no se podían subir soportes
al corte; los demás tienen que ver con buscar cortes, escribir actividades propias y detalles
de la pantalla. Los siete se corrigieron en el prototipo y cada uno queda explicado con su cambio y sus
capturas de antes y después.

La mayor parte de los comentarios no fueron fallas, sino necesidades del trabajo diario que
el prototipo aún no cubre: herramientas y equipos, materiales por lote, permisos por obra y
cotizaciones. Quedan anotadas como pedidos para próximas versiones, y refinamiento del entendimiento del funcionamiento de loard.

\end{document}
